\chapter{\label{I_3}Le projet d'automatiser l'indexation des archives audiovisuelles}

%I 3 

L'architecture modélisée par le \gls{FRBR} décrite jusqu'ici conditionne la structure de la description, et de fait, de l'indexation de cette description selon des normes définies. En effet, le modèle de données permet d'organiser l'indexation en couches intelligibles, permettant à l'utilisateur de se repérer au sein d'un logiciel de gestion des collections, et de faciliter les entrées de description. Les tâches d'indexation peuvent toutefois s’avérer complexe et chronophage, spécifiquement dans le cas de corpus audiovisuels tels que celui que nous avons analysé, poussant l'éditeur à émettre le projet de développement d'outils automatisés permettant d'assister l'utilisateur.



  
\section{\label{I_3_A}Indexation} 

L’indexation documentaire peut être définie de la manière suivante :  

\begin{quote} 
	
	"L'indexation est cette technique consistant à caractériser le contenu d'un document et l'information qu'il détient de manière à le retrouver quand on effectue des recherches sur l'un des sujets dont il traite. [...] L'indexation permet de s'orienter dans la masse des documents et d'organiser ses connaissances. "\footcite{zotero-item-582} 
	
\end{quote} 

Afin de "fournir un point d'accès sûr et vérifié à l'information", l'indexation implique la mise en relation des termes produits par la description avec des référentiels, des thésaurus et des listes d'autorité, afin de pouvoir les intégrer au modèle de données.\footcite{zotero-item-675} Ainsi, par exemple, le réalisateur d’un film est indexé au niveau Œuvre, par une liste contrôlée d’autorités prévue par l’outil de gestion des collections.  

La structuration des métadonnées d’indexation permet d'améliorer la découvrabilité de ces métadonnées via des outils de recherche et la navigation entre les différentes connaissances produites.\footcite{zotero-item-675} Une indexation reliée à des thésaurus contrôlés permet notamment à l’utilisateur d’effectuer des requêtes croisées complexes au sein de l’outil de gestion des collections. Cette tâche, prévue par le logiciel, effectue le lien indispensable entre la description produite et l’exploitabilité des données par l’utilisateur via des outils de recherche. À terme, l'indexation permet l’interopérabilité des données et leur partage sur le web. \footcite{zotero-item-675}  

Toutefois, la nature des archives audiovisuelles complexifie l’indexation pour l’utilisateur. En effet, nous l’avons évoqué, de telles archives se constituent comme des “boites noires”, nécessitant le matériel de lecture adéquat afin d’en extraire le sens et de pouvoir l’indexer. De plus, l’indexation du corpus que nous avons décrit, composé de rushes documentaires, suppose en amont la division du flux vidéo en différentes unités de sens thématiques. La description ne peut pas s’appliquer de manière précise dans le cas d’un document vidéo brut contenant de multiples séquences spatio-temporelles. Dans le cas de ce type d’archives audiovisuelles, la segmentation temporelle constituerait alors la condition préalable à toute opération d’indexation.
\section{\label{I_3_B}Segmentation} 

Dans le cas des archives audiovisuelles, une indexation fine ne peut s’effectuer à l’échelle du fichier brut dans sa totalité, en particulier lorsqu’il s’agit de rushes contenant une succession importante de séquences non montées. Afin de rendre le contenu de telles archives exploitable au sein d’un modèle de données cohérent, il est indispensable d’analyser finement le flux continu de l’image via une tâche de segmentation. Cette opération est à l’origine de questionnements autour de l’intelligibilité du document audiovisuel :  

\begin{quote} 
	
	"Quels sont les mécanismes qui permettent d’in-terpréter cette longue séquence d’images [...] en une séquence d’actions dramatiques, reliées entre elles par des liens de causalité, et qui nous permettent d’en comprendre le sens (l’histoire) ? "\footcite{zotero-item-644} 
	
\end{quote} 

L’intelligibilité du document audiovisuel passerait ainsi par une interprétation, via une segmentation de l’enchaînement des images afin d’en extraire le sens. La segmentation revêt, dans le cadre de l’indexation d’archives audiovisuelles, une dimension essentiellement temporelle, dans l'optique de découper un flux vidéo selon certains critères. Il s'agit d'un découpage conceptuel et non d'un montage technique qui modifie le sens de l’image animée. La segmentation s'imposerait comme la condition d'une indexation fine puisqu'en apposant à des séquences des balises temporelles, elle permettrait de rendre possible "la production de descriptions structurées". \footcite{stockinger2011} Ces balises temporelles délimitent des ensembles homogènes au sein d'une archive audiovisuelle, à la manière de chapitres dans un ouvrage textuel \footcite{stockinger2011}. Si la structuration de certains types de contenus audiovisuels tels que les journaux télévisés, le sport, la publicité, sont "relativement stable et explicite de leur structure", et peuvent alors facilement être segmentés, les archives audiovisuelles regroupent des corpus divers et variés à la structure aléatoire qui rendent complexe cette tâche \footcite{carrive2014} 

La segmentation s’applique particulièrement bien au corpus d’archives audiovisuelles composé de rushes documentaires, que nous avons décrit. En effet, au sein du logiciel prévu par Skinsoft, la segmentation est pratiquée par l’institution patrimoniale préservant ces archives, qui associe des repères temporels à des extraits numérisés et issus de rushes filmés plus larges. Cette pratique permet d’extraire des rushes certaines séquences jugées importantes à indexer de manière précise. Ces séquences sont ensuite rendues intelligibles par l’indexation de leur contenu précis et deviennent requêtables de manière autonome au sein du modèle de données.  

La segmentation répond donc à un objectif fixé par l’indexation, celui d’agir comme “un filtre de pertinence" sur les éléments à indexer \footcite{stockinger2011}. Sa mise en place suppose d’en définir la granularité (analyse au niveau du plan, du changement de plan, de scènes, de séquences, etc), afin d’isoler les éléments d’analyse \footcite{stockinger2011}. Ce travail d’analyse s’articule autour de trois types de segmentation, conditionnée par la nature du médium cinématographique. En effet, la segmentation temporelle est non seulement conditionnée par l’image mais également par le son, composant essentiel des documents audiovisuels \footcite{carrive2014}. La segmentation temporelle suppose également une segmentation spatiale, induisant l'analyse de la situation des éléments dans l'image \footcite{zhou2022}. Enfin, une segmentation sémantique intervient, apposant une description aux éléments extraits, permettant de produire un découpage logique \footcite{zhou2022}. Les catégories sémantiques sont prédéfinies par l'utilisateur et peuvent être de plusieurs types conformément à la nature de l’image animée. Ces catégories peuvent être d’ordre technique, relationnelles, de contenu, de mouvement, etc.  

Parmi ces trois types de segmentation, seule la segmentation temporelle est spécifique à l'image animée, puisqu'elle analyse la différence temporelle et de mouvement entre les plans ainsi que les relations qu’ils entretiennent, afin de pouvoir en extraire des séquences distinctes. Les segmentation spatiales et sémantiques sont, en revanche, également applicables à l'image fixe. 

Ce processus de segmentation s’inscrit directement dans l’architecture du modèle FRBR. Les segments temporels créés sont décrits et indexés comme des nouvelles entités et rattachés aux entités pertinentes du modèle. C’est dans le contexte de l’indexation de grands volumes d’archives et spécifiquement de corpus complexes que peut être envisagée la mise en place d’une automatisation de la segmentation, puis de l‘indexation. Selon Attard, la mise en place d’une segmentation automatisée permettrait “d’accélérer le traitement de grands volumes vidéo tout en garantissant la cohérence du processus” (traduction : automated segmentation accelerates the processing of large volumes of video content and ensures consistency and accuracy in the segmentation process." ) \footcite{attard2025}. 

Toutefois, l’efficacité promise par l’automatisation de la segmentation et de l’indexation ne pourra être mise en place qu’à la condition selon laquelle les algorithmes d’apprentissage automatique puissent produire ces tâches selon des règles métiers précises et un modèle de données explicite. En tant qu’éditeur de logiciel de gestion des collections, Skinsoft ambitionne de résoudre ces complexités afin d’intégrer ces nouvelles fonctionnalités au sein de l’espace de travail S-Movies. Dès lors, il convient d’examiner les directions adoptées par l’éditeur afin de proposer une automatisation innovante conditionnée par une structure de données complexe.  
\section{\label{I_3_C}Le projet \gls{ia} chez Skinsoft} 

\subsection{Un projet global}
La mise en place d'une indexation automatique, visant à faciliter le travail des utilisateurs métier et à le rendre efficace, s'inscrit dans un projet de mise en place de l'intelligence artificielle à une plus large échelle au sein de l'application chez Skinsoft, afin d'automatiser et de faciliter les pratiques de gestion des collections. L'intelligence artificielle (\gls{ia}) peut être définie ainsi : 

\begin{quote} 
	
	Par IA, on entend, dans une acception large et ouverte, la tentative de construire des machines et des systèmes capables de penser et d'agir rationnellement \footcite{zotero-item-584}. 
	
\end{quote} 



L'objectif serait alors, pour un éditeur comme Skinsoft, celui de rendre ces systèmes capables de pouvoir raisonner afin d'effectuer les mêmes actions que l'homme. Dans le cadre d'un système de gestion des collections, cela signifie que les systèmes d'intelligence artificielle devront effectuer des tâches documentaires complexes telles que nous les avons décrites. Cette implémentation doit donc s'inscrire dans cadre défini donnant la direction de développement de ces machines, adapté au contexte patrimonial. 

L'implémentation de l'\gls{ia} au sein de l'application n'est pas seulement issue d'une logique commerciale d'innovation, elle est centrée avant tout sur l'amélioration de l'expérience utilisateur, en vue de rendre l'indexation des collections efficace. L'une des caractéristiques majeures de l'application, nous l'avons évoqué, est sa flexibilité et son adaptabilité aux spécificités des collections et à diverses règles métier. La mise en place de l'\gls{ia} permettrait de mettre cette adaptabilité en place plus efficacement et à plus large échelle. Automatiser la migration de données, le paramétrage de l'application, la navigation dans l'application permettrait à l'entreprise de satisfaire ses clients plus rapidement et de gérer ainsi de multiples contrats. Cette volonté d'implémenter l'\gls{ia} s'inscrit notamment dans l'appel à projets "France 2030" qui couvre la maitrise des technologies numériques et le "déploiement de l'innovation" \footcite{zotero-item-636}. Ce projet finance des initiatives à l'intersection de la culture, du numérique et de l'\gls{ia}. 

\subsection{Identifier le besoin pour instaurer un cadre éthique}

Toutefois, la question du besoin est centrale, pour une utilisation de l'\gls{ia} éthique. Comme le mentionne le Ministère de la culture : "il ne s’agit pas d’utiliser l’IA par défaut, mais d’abord d’examiner le besoin à couvrir" \footcite{2025a}. L'objectif est double pour Skinsoft. La mise en place de l'\GLS{ia} dans l'application permettrait d'une part d'autonomiser les institutions sur leur usage du logiciel et d'autre part d'améliorer l'efficacité de l'indexation de leurs collections. Ainsi, les équipes Skinsoft seraient moins sollicitées pour des difficultés liées à l'usage de l'application qui pourraient être réglés par une machine intelligente. Afin de répondre à ces objectifs, Skinsoft prévoie notamment la mise en place d'une recherche en langage naturel, c'est-à-dire de permettre à l'utilisateur d'effectuer des recherches plus facilement en utilisant un langage naturel---son propre langage--- plutôt qu'un langage contrôlé. Cette fonctionnalité est abordée dans la spécification fonctionnelle intégrée à l'annexe A (\ref{annexe_A}). L'objectif est ainsi de rendre l'utilisation globale de l'application plus simple, et d'améliorer l'efficacité de la gestion de collections. 

Toutefois, l'implémentation de l'\gls{ia} pour la création de contenus structurés soulève des enjeux éthiques dès lors que l'objectif de l'indexation est de produire des connaissances pour la transmission auprès du public. De plus, les pratiques documentaires, nous l'avons vu, sont des pratiques normalisées, structurées par un modèle de données afin d'assurer leur intéropérabilité.  

Ainsi, Skinsoft porte des points d'attention au projet. Les moteurs d'intelligence artificielle implémentés devraient notamment faciliter le travail des utilisateurs et être construits comme des outils de proposition, sans réaliser entièrement eux-mêmes la production de connaissance. L'humain est ainsi toujours intégré au sein de la chaîne de travail dans la définition du projet, en étant le réel décisionnaire final de ce que produit l'\gls{ia}.  

\subsection{Le besoin d'une segmentation et d'une indexation automatisées}

Concrètement, Skinsoft souhaite mettre en place différents moteurs d'intelligence artificielle, entraînés sur des données patrimoniales, et dédiés à l'automatisation partielle de tâches métier dans l'application. Parmi ces tâches, il est prévu l'automatisation de l'indexation au sens large au sein de l'application, et notamment appliquée aux archives audiovisuelles, dans S-Movies. Si, pour les images fixes, peut être mise en place une analyse automatique de leur contenu afin de l'indexer, l'indexation automatique des archives audiovisuelles se découpe,en théorie, en plusieurs sous-tâches, propres à la nature du médium cinématographique.  

Ainsi, le moteur d'intelligence artificielle dédié à l'indexation des archives audiovisuelles au sens large, peut être divisé en plusieurs sous-moteurs, tous dédiés à l'analyse de la coordination du son avec le mouvement des images et leur contenu. Concrètement, cela peut désigner les tâches automatiques suivantes : transcription et analyse du son, reconnaissance des locuteurs, segmentation et analyse des images en mouvement, la transcription du texte présent dans l'image. Les données extraites de ces différentes tâches devront être automatiquement inscrites au sein de référentiels afin d'assurer leur interopérabilité. Par exemple, lors de la reconnaissance d'un locuteur, son nom doit pouvoir être relié à sa notice d'autorité dans l'application.  

Toutefois, le projet d'automatiser l'indexation est issu d'une demande spécifique d'une institution cliente du logiciel, gestionnaire du corpus audiovisuel spécifique que nous avons décrit, constitué de \gls{rush}, nécessitant également un processus de segmentation. Or, la mise en place de la segmentation suppose d’en définir la granularité d'analyse au niveau de la syntaxe audiovisuelle, composée d'éléments divers : le \gls{plan}, la \gls{sequence}, la scène, etc. \footcite{stockinger2011}. Dans le cadre de ce corpus, la segmentation manuelle s'effectue déjà au niveau de la \gls{sequence}. Une \gls{sequence}, en termes cinématographiques, peut être définie comme "une unité d'action, un fragment du film qui raconte en plusieurs plans une suite d'événements isolable dans la construction narrative." \footcite{journot2008}. Dans ce contexte, la segmentation désigne donc la division d'images en mouvements en plusieurs \gls{sequence}s. 

L'institution souhaiterait ainsi avoir à sa disposition un outil qui détecterait automatiquement les changements de \gls{sequence}s, pour pouvoir les décrire individuellement de manière fine, évitant ainsi une description trop globale au niveau du \gls{rush}. Comme nous l'avons évoqué, des marqueurs temporels sont déjà présents dans les données de l'institution, au sein des notices Oeuvres, afin de délimiter les \gls{sequence}s pertinentes. L'objectif est alors celui d'automatiser l'apposition de marqueurs temporels sur la vidéo en détectant les changements majeurs d'action, de lieu, d'objets ou de personnages, marquant le passage d'une séquence à une autre. 

À ce besoin s'ajoute celui d'indexer automatiquement les \gls{sequence}s identifiées par des mots-clés, afin de leur donner une valeur sémantique. En effet, les \gls{sequence}s, au sein de la syntaxe cinématographique, désignent une unité à la fois temporelle et sémantique, marquant un changement de forme significatif au sein du flux audiovisuel, mais également un changement de fond et de sens donné à l'image. Ainsi, la segmentation ne suffit pas à les qualifier de \gls{sequence}s, puisque c'est véritablement leur description qui leur redonne leur sens. Cette description doit ainsi pouvoir être automatisée par la production de mots-clés, puis indexée en lien avec les normes métier employées.

Ainsi, le besoin d'automatisation identifié se divise en deux tâches spécifiques : la segmentation du flux audiovisuel en \gls{sequence}s, et leur indexation par mots-clés. L'automatisation de ces tâches permettrait une accélération du traitement des fonds audiovisuels, et à terme, une mise à disposition plus large auprès d'utilisateurs. Si le projet d'automatiser ces tâches est issu d'un besoin spécifique à uen institution, il pourra être généralisé à toute institution qui préserve des collections audiovisuelles. 



Le traitement des collections audiovisuelles exige une articulation entre la segmentation temporelle d'une part, et l'indexation des \gls{sequence}s produites, d'autre part. En effet, la segmentation s'impose comme la condition préalable indispensable à une description documentaire exploitable par les utilisateurs. À travers la construction d'un cadre pour le projet d'implémentation de l'\gls{ia} au sein de son logiciel, Skinsoft souhaite concevoir, à partir des besoins métier, un outil d'assistance pour le traitement massifs des fonds audiovisuels. L'identification des besoins d'une part, et du cadre théorique d'implémentation d'autre part, permet d'imaginer un processus automatique qui s'inscrit au sein de la structure documentaire du \gls{FRBR}. 


